\documentclass[10pt]{article}
\usepackage[T1]{fontenc}
\usepackage{lmodern,amsmath,amssymb,amsthm,microtype}
\usepackage[margin=0.85in]{geometry}
\usepackage[colorlinks=true,urlcolor=blue]{hyperref}
\newtheorem{theorem}{Theorem}
\title{Even-Block Partitions: Modular Periodicity and Bessel Moments}
\author{Qichao Wang\\{\small Hebei University of Technology}\\[0.3em]
Daoyu Dong\\{\small University of Electronic Science and Technology of China}}
\date{October 26, 2026}
\begin{document}
\maketitle
\begin{abstract}
For the number of partitions of a $2n$-element set into even blocks, we prove
eventual periodicity modulo every positive integer by an explicit finite
differential module. We also prove the proposed Bessel moment formula and
the stated six-term eventual period modulo 10. These resolve the two
conjectures recorded in OEIS A005046.
\end{abstract}

\section{A finite differential module}
Let $a(n)$ count partitions of $\{1,\ldots,2n\}$ into even blocks, with
$a(0)=1$. Choosing the block containing 1 gives the exact recurrence
\[
 a(n)=\sum_{j=1}^n\binom{2n-1}{2j-1}a(n-j).
\]
The exponential formula for set partitions, or this recurrence, gives
\[
 F(x)=\exp(\cosh x-1)=\sum_{n\ge0}a(n)\frac{x^{2n}}{(2n)!}.
\]
Define integer polynomials $P_0(s,c)=1$ and
\[
 P_{r+1}=L P_r,\qquad L=c\partial_s+s\partial_c+s.
\]
Here the final $s$ denotes multiplication. The chain and product rules give
\[
 F^{(r)}(x)=F(x)P_r(\sinh x,\cosh x),\qquad a(n)=P_{2n}(0,1).
\]

\begin{theorem}
For every integer $m\ge1$, the sequence $a(n)\bmod m$ is eventually periodic.
For $m\ge2$ its values are obtained by iteration of one explicitly defined
linear map on a module with $m^{m^2}$ elements. In particular, a preperiod
$u\ge0$ and period $v\ge1$ exist with $u+v\le m^{m^2}$.
\end{theorem}
\begin{proof}
Work over $R=\mathbb Z/m\mathbb Z$ and form the finite free $R$-module
\[
 V_m=R[s,c]/(s^m,c^m-1).
\]
Its basis consists of the $m^2$ monomials $s^ic^j$ with $0\le i,j<m$.
The derivation $\delta=c\partial_s+s\partial_c$ preserves the ideal:
\[
 \delta(s^m)=m s^{m-1}c=0,\qquad
 \delta(c^m-1)=m c^{m-1}s=0.
\]
Multiplication by $s$ also preserves it, so $L$ and $B=L^2$ are well-defined
linear maps on $V_m$. Evaluation $\varepsilon(P)=P(0,1)$ is well-defined
because both ideal generators vanish there. Thus
\[
 a(n)\equiv\varepsilon(B^n1)\pmod m.
\]
There are only $m^{m^2}$ possible states. Among the first $m^{m^2}+1$
states two coincide; applying $B$ repeatedly propagates the equality
forever. Their indices give $u,v$ and the stated bound. For $m=1$ the
sequence is identically zero. This proves Bala's 2023 conjecture~\cite{oeis}.
\end{proof}
The construction is an exact finite algorithm for arbitrary $m$, although its
state bound is deliberately coarse. It does not assume a prime modulus or
divide by factorials modulo $m$.

\section{The period modulo ten}
The small quotient rings yield short certificates for the example in the entry.
In $V_2$ direct application of $L$ gives
\[
 P_0=1,\quad P_2=c,\quad P_4=1+c,\quad P_6=1.
\]
Thus $a(n+3)=a(n)\pmod2$ for every $n\ge0$.
In $V_5$ the corresponding calculation gives
\[
 P_2=c+s^2,\qquad
 P_4=c+3c^2+4s^2+s^2c+s^4,\qquad P_6=c+s^2.
\]
Consequently $a(n+2)=a(n)\pmod5$ for $n\ge1$. The Chinese remainder
theorem and the first six positive-index values now give
\[
 a(n)\bmod10=(1,4,1,9,6,9)_{\,1+((n-1)\bmod6)},\qquad n\ge1.
\]
Its least eventual period is 6: the displayed block has no period 1, 2 or 3.

\section{The Bessel moment formula}
Define the modified Bessel values, for integers $k\ge0$, by their absolutely
convergent series
\[
 I_k(1)=\sum_{j\ge0}\frac{1}{2^{2j+k}j!(j+k)!}.
\]
\begin{theorem}
For every integer $n\ge1$,
\[
 a(n)=2e^{-1}\sum_{k=1}^{\infty} k^{2n}I_k(1).
\]
At $n=0$ the corresponding identity is
$1=e^{-1}\bigl(I_0(1)+2\sum_{k\ge1}I_k(1)\bigr)$.
\end{theorem}
\begin{proof}
Multiply the absolutely convergent series for $\exp(e^x/2)$ and
$\exp(e^{-x}/2)$ and collect terms with difference of indices equal to $k$.
This gives the Laurent expansion
\[
 e^{\cosh x}=\sum_{k\in\mathbb Z} I_{|k|}(1)e^{kx}.
\]
All regroupings and derivatives are valid locally uniformly for complex $x$.
Indeed, on $|x|\le R$, the absolute value of an $r$-th derivative of the
double-series term with indices $p,q$ is at most
\[
 \frac{(p+q)^r e^{R(p+q)}}{2^{p+q}p!q!},
\]
whose sum is $\sum_{t\ge0}t^r e^{Rt}/t!<\infty$
(with the usual constant-term interpretation when $r=0$).
Differentiating $2n$ times at zero and multiplying by $e^{-1}$ gives the
claim, since the $k=0$ term vanishes for $n>0$ and opposite indices agree.
This proves Yanev's 2026 conjecture~\cite{oeis}.
\end{proof}

\section{Verification and scope}
The standard-library verifier compares the partition recurrence with a stored
132-term official prefix and with literal set partitions through ten elements.
It independently implements the finite differential modules modulo 2 through
32 through index 250, verifies the two whole-polynomial cycle certificates,
and checks the Bessel formula numerically for $1\le n\le12$.
Numerical agreement is a consistency check; convergence and equality are proved
above. The formulas are attributed to their OEIS proposers. No claim of an
exhaustive literature-priority search or of OEIS acceptance is made.
\begin{thebibliography}{9}
\bibitem{oeis} The OEIS Foundation, A005046, with conjectures by
P. Bala (April 12, 2023) and V. Yanev (February 9, 2026),
\url{https://oeis.org/A005046}, accessed October 26, 2026.
\end{thebibliography}
\end{document}
